\documentclass[12pt,a4paper]{article}
\usepackage[utf8]{inputenc}
\usepackage[margin=1in]{geometry}
\usepackage{listings}
\usepackage{xcolor}
\usepackage{booktabs}
\usepackage{array}
\usepackage{tabularx}
\usepackage{fancyhdr}
\usepackage{graphicx}
\usepackage{amsmath}

% Code styling
\definecolor{codegreen}{rgb}{0,0.6,0}
\definecolor{codegray}{rgb}{0.5,0.5,0.5}
\definecolor{codepurple}{rgb}{0.58,0,0.82}
\definecolor{backcolour}{rgb}{0.95,0.95,0.95}

\lstdefinestyle{pythonstyle}{
    backgroundcolor=\color{backcolour},
    commentstyle=\color{codegreen},
    keywordstyle=\color{blue}\bfseries,
    numberstyle=\tiny\color{codegray},
    stringstyle=\color{codepurple},
    basicstyle=\ttfamily\footnotesize,
    breakatwhitespace=false,
    breaklines=true,
    captionpos=b,
    keepspaces=true,
    numbers=left,
    numbersep=5pt,
    showspaces=false,
    showstringspaces=false,
    showtabs=false,
    tabsize=2,
    frame=single,
    rulecolor=\color{gray}
}

\lstset{style=pythonstyle}

\pagestyle{fancy}
\fancyhf{}
\fancyhead[C]{\textbf{SHAFE:JAWAD}}
\fancyfoot[C]{Page \thepage}
\renewcommand{\headrulewidth}{0.4pt}

\begin{document}

\begin{center}
{\LARGE \textbf{Keras Output Layers Cheat Sheet}}\\[0.3em]
{\large Linear vs Logistic vs Multi-class}
\end{center}

\vspace{0.5cm}

%----------------------------------------------------------------------
% PAGE 1: LINEAR REGRESSION
%----------------------------------------------------------------------

\section{Linear Regression}

Predicts continuous values (like prices, temperatures). Output can be any real number, so we use \textbf{linear activation} (basically no activation).

\subsection{Output Layer}

\begin{lstlisting}[language=Python, caption=Linear Regression Output Layer]
from tensorflow.keras import layers

# Output layer for regression
output_layer = layers.Dense(
    units=1,              # Single output neuron
    activation='linear'   # No activation (unbounded output)
)
\end{lstlisting}

\subsection{Compilation}

\begin{lstlisting}[language=Python, caption=Compiling for Regression]
model.compile(
    optimizer='adam',
    loss='mse',           # Mean Squared Error
    metrics=['mae']       # Mean Absolute Error
)
\end{lstlisting}

\subsection{Quick Facts}

\begin{itemize}
    \item \textbf{Output neurons:} 1
    \item \textbf{Activation:} Linear $\rightarrow f(x) = x$
    \item \textbf{Loss:} MSE or MAE
    \item \textbf{Output range:} $(-\infty, +\infty)$
\end{itemize}

\subsection{What if output is constrained?}

\begin{lstlisting}[language=Python, caption=Constrained Regression Outputs]
# If target >= 0 (e.g., prices, counts)
layers.Dense(1, activation='relu')

# If target in (0, 1) (e.g., probabilities, ratios)
layers.Dense(1, activation='sigmoid')

# If target in (-1, 1)
layers.Dense(1, activation='tanh')
\end{lstlisting}

\newpage

%----------------------------------------------------------------------
% PAGE 2: LOGISTIC REGRESSION (BINARY CLASSIFICATION)
%----------------------------------------------------------------------

\section{Logistic Regression (Binary)}

Two classes only (yes/no, spam/not spam). Use \textbf{sigmoid} to squash output into probability between 0 and 1.

\subsection{Output Layer}

\begin{lstlisting}[language=Python, caption=Binary Classification Output Layer]
from tensorflow.keras import layers

# Output layer for binary classification
output_layer = layers.Dense(
    units=1,              # Single output neuron
    activation='sigmoid'  # Squashes output to (0, 1)
)
\end{lstlisting}

\subsection{Compilation}

\begin{lstlisting}[language=Python, caption=Compiling for Binary Classification]
model.compile(
    optimizer='adam',
    loss='binary_crossentropy',  # Log loss for binary
    metrics=['accuracy']
)
\end{lstlisting}

\subsection{Quick Facts}

\begin{itemize}
    \item \textbf{Output neurons:} 1
    \item \textbf{Activation:} Sigmoid $\rightarrow \sigma(x) = \frac{1}{1 + e^{-x}}$
    \item \textbf{Loss:} \texttt{binary\_crossentropy}
    \item \textbf{Output range:} $(0, 1)$ = probability of class 1
\end{itemize}

\subsection{Getting Predictions}

\begin{lstlisting}[language=Python, caption=Binary Classification Predictions]
# Get probability
prob = model.predict(X_test)  # Output: [[0.87], [0.12], ...]

# Convert to class labels
predictions = (prob > 0.5).astype(int)

# Custom threshold
predictions = (prob > 0.7).astype(int)
\end{lstlisting}

\subsection{Labels}

\begin{lstlisting}[language=Python, caption=Binary Label Format]
# Labels are just 0 or 1
y_train = [0, 1, 1, 0, 1, 0, ...]  # Shape: (n_samples,)
\end{lstlisting}

\newpage

%----------------------------------------------------------------------
% PAGE 3: MULTI-CLASS CLASSIFICATION
%----------------------------------------------------------------------

\section{Multi-class Classification}

More than 2 classes (digits 0-9, cat/dog/bird, etc). Use \textbf{softmax} to get a probability for each class (they all sum to 1).

\subsection{Output Layer}

\begin{lstlisting}[language=Python, caption=Multi-class Classification Output Layer]
from tensorflow.keras import layers

num_classes = 10  # e.g., digits 0-9

# Output layer for multi-class
output_layer = layers.Dense(
    units=num_classes,    # One neuron per class
    activation='softmax'  # Outputs sum to 1
)
\end{lstlisting}

\subsection{Compilation}

\begin{lstlisting}[language=Python, caption=Compiling for Multi-class Classification]
# If labels are integers (0, 1, 2, ...) - most common!
model.compile(
    optimizer='adam',
    loss='sparse_categorical_crossentropy',
    metrics=['accuracy']
)

# If labels are one-hot encoded
model.compile(
    optimizer='adam',
    loss='categorical_crossentropy',
    metrics=['accuracy']
)
\end{lstlisting}

\subsection{Quick Facts}

\begin{itemize}
    \item \textbf{Output neurons:} K (one per class)
    \item \textbf{Activation:} Softmax $\rightarrow \frac{e^{z_i}}{\sum_{j} e^{z_j}}$
    \item \textbf{Loss:} \texttt{sparse\_categorical\_crossentropy} (int) /\\ \texttt{categorical\_crossentropy} (one-hot)
    \item \textbf{Output:} K probabilities that sum to 1
\end{itemize}

\subsection{Labels (two options)}

\begin{lstlisting}[language=Python, caption=Multi-class Label Formats]
# Integer labels (use sparse_categorical_crossentropy)
y_train = [0, 2, 1, 5, 3, ...]  # Shape: (n_samples,)

# One-hot labels (use categorical_crossentropy)
from tensorflow.keras.utils import to_categorical
y_train_onehot = to_categorical(y_train, num_classes=10)
# Shape: (n_samples, num_classes)
# Example: class 2 -> [0, 0, 1, 0, 0, 0, 0, 0, 0, 0]
\end{lstlisting}

\newpage

%----------------------------------------------------------------------
% PAGE 4: COMPARISON TABLE
%----------------------------------------------------------------------

\section{Side-by-Side Comparison}

\subsection{The Big Table}

\begin{table}[h!]
\centering
\small
\renewcommand{\arraystretch}{1.5}
\newcolumntype{Y}{>{\raggedright\arraybackslash}X}
\begin{tabularx}{\textwidth}{|l|Y|Y|Y|}
\hline
\textbf{Aspect} & \textbf{Linear Regression} & \textbf{Logistic Regression} & \textbf{Multi-class} \\
\hline
\textbf{Task} & Predict continuous value & Binary classification & K-class classification \\
\hline
\textbf{Output Neurons} & 1 & 1 & K (num classes) \\
\hline
\textbf{Activation} & \texttt{linear} & \texttt{sigmoid} & \texttt{softmax} \\
\hline
\textbf{Loss Function} & \texttt{mse} / \texttt{mae} & \texttt{binary\_\allowbreak crossentropy} & \texttt{categorical\_\allowbreak crossentropy} \\
\hline
\textbf{Output Range} & $(-\infty, +\infty)$ & $(0, 1)$ & $(0, 1)$ each, sum $= 1$ \\
\hline
\textbf{Label Format} & Continuous values & 0 or 1 & Integers or one-hot \\
\hline
\textbf{Example Use} & House prices, temperature & Spam detection, medical diagnosis & Image classification, sentiment \\
\hline
\end{tabularx}
\caption{Comparison of Neural Network Output Configurations}
\end{table}

\subsection{Activation Functions}

\begin{table}[h!]
\centering
\renewcommand{\arraystretch}{1.5}
\begin{tabular}{|l|l|l|}
\hline
\textbf{Function} & \textbf{Formula} & \textbf{Range} \\
\hline
Linear & $f(x) = x$ & $(-\infty, +\infty)$ \\
\hline
Sigmoid & $\sigma(x) = \frac{1}{1 + e^{-x}}$ & $(0, 1)$ \\
\hline
Softmax & $\frac{e^{z_i}}{\sum_j e^{z_j}}$ & $(0, 1)$, sum = 1 \\
\hline
ReLU & $f(x) = \max(0, x)$ & $[0, +\infty)$ \\
\hline
Tanh & $\tanh(x) = \frac{e^x - e^{-x}}{e^x + e^{-x}}$ & $(-1, 1)$ \\
\hline
\end{tabular}
\caption{Common Activation Functions}
\end{table}

\subsection{How to Pick?}

\begin{enumerate}
    \item \textbf{Continuous target?} (price, temp, etc.) $\rightarrow$ \textbf{Linear}
    \begin{itemize}
        \item \texttt{Dense(1, activation='linear')} + \texttt{loss='mse'}
    \end{itemize}

    \item \textbf{2 classes?} (yes/no, 0/1) $\rightarrow$ \textbf{Sigmoid}
    \begin{itemize}
        \item \texttt{Dense(1, activation='sigmoid')} + \texttt{loss='binary\_crossentropy'}
    \end{itemize}

    \item \textbf{K classes?} (K > 2) $\rightarrow$ \textbf{Softmax}
    \begin{itemize}
        \item \texttt{Dense(K, activation='softmax')}\\
              \texttt{loss='sparse\_categorical\_crossentropy'}
    \end{itemize}
\end{enumerate}

\vfill

\noindent\fbox{\parbox{\dimexpr\textwidth-2\fboxsep-2\fboxrule}{%
\textbf{Remember:} Hidden layers $\rightarrow$ use ReLU or tanh (your choice).\\
Output layer $\rightarrow$ depends on the task (linear/sigmoid/softmax)!
}}

\end{document}